% ---------------------------------------------------------------------------
%  T1 — Actividad 1 (Ética, parte 1)
%  Presentación: "Ética en la Publicación Científica"
%  Grupo 1 de Cátedra — Metodologías de Investigación Aplicada (MIA)
%  Tema: usach (normas gráficas MNG 2024-2S)
%  Compilar: latexmk -pdf -outdir=build main.tex
% ---------------------------------------------------------------------------
\documentclass[11pt,aspectratio=169,xcolor=table]{beamer}

\IfFileExists{spanish.ldf}{%
  \usepackage[spanish,es-noquoting,es-nodecimaldot]{babel}%
}{%
  \usepackage[T1]{fontenc}
  \renewcommand{\figurename}{Figura}
  \renewcommand{\tablename}{Tabla}
}
\usepackage[utf8]{inputenc}
\usepackage{amsmath,amssymb}
\usepackage{booktabs,tabularx}
\usepackage{microtype}

\usetheme{usach}

% --- Metadatos --------------------------------------------------------------
\title[Ética en la publicación científica]{Ética en la Publicación Científica}
\subtitle{De la manipulación individual de datos a la industrialización del fraude en informática}
\author[G.~Ahumada · D.~Fernández · I.~Fernández · P.~Figueroa]{%
  Gonzalo Ahumada · Diego Fernández · Iván Fernández · Pablo Figueroa}
\date{Septiembre de 2026}

\usachfacultad{Facultad de Ingeniería}
\usachdepartamento{Departamento de Ingeniería Informática}
\usachprograma{Magíster en Ingeniería Informática}
\usachasignatura{Metodologías de Investigación Aplicada}
\usachprofesor{Prof. Roberto Ibáñez González}

\institute[USACH]{Universidad de Santiago de Chile}

\begin{document}

\usachportada
\note{Buenos días. Somos el Grupo 1 de cátedra: Diego Fernández, Iván
Fernández, Pablo Figueroa y quien les habla. Presentamos la primera
evaluación de ética: dos casos reales centrados en la publicación como
producto — el histórico de Jan Hendrik Schön y el contemporáneo de
retracciones masivas en la ACM.}

% --- Agenda ------------------------------------------------------------------
\begin{frame}{Contenido}
  \tableofcontents
\end{frame}

% =============================================================================
\section{Marco normativo}
% =============================================================================

\begin{frame}{La publicación como producto de la investigación}{Tríada FFP y delimitación del análisis}
  \begin{block}{Foco del análisis}
    La publicación científica es el producto formal y acreditable de la
    investigación: en ella se materializan la validez, la novedad y el rigor
    de los hallazgos.
  \end{block}
  \begin{itemize}\usachaire
    \item \textbf{Fabricación}: invención de datos o resultados inexistentes.
    \item \textbf{Falsificación}: manipulación de procesos o eliminación de datos reales.
    \item \textbf{Plagio}: apropiación de ideas, texto o código sin atribución.
  \end{itemize}
  \alert{Delimitación}: el foco es la \emph{publicación como producto}; se
  excluyen las faltas contra seres vivos y humanos (segunda actividad).
  \note{Basta con dos documentos clave: \emph{On Being a Scientist} y la
  Declaración de Singapur. La tríada FFP resume la mala conducta clásica y
  delimita nuestro análisis a la publicación como producto.}
\end{frame}

\begin{frame}{Deberes de Singapur y nuevos fraudes}{Del fraude artesanal al fraude industrializado}
  \begin{columns}[T,onlytextwidth]
    \begin{column}{0.50\textwidth}
      \textbf{Declaración de Singapur (deberes clave)}
      \begin{itemize}
        \item Resp.~4: custodia transparente de datos primarios.
        \item Resp.~6: autoría con responsabilidad plena.
        \item Resp.~8: arbitraje imparcial y riguroso.
      \end{itemize}
    \end{column}
    \begin{column}{0.46\textwidth}
      \textbf{Presiones y nuevas modalidades}
      \begin{itemize}
        \item \emph{Publish or perish}: el volumen como incentivo perverso.
        \item Fraude \alert{artesanal}: el científico que falsifica en soledad.
        \item Fraude \alert{industrializado}: \emph{paper mills} que venden autorías.
        \item Vigilancia comunitaria: Retraction Watch y PubPeer.
      \end{itemize}
    \end{column}
  \end{columns}
  \note{Singapur exige custodia de datos, autoría solidaria y revisión por
  pares rigurosa. Distinguimos el fraude artesanal clásico del fraude
  industrializado que domina hoy la informática.}
\end{frame}

% =============================================================================
\section{Caso 1: Jan Hendrik Schön}
% =============================================================================

\begin{frame}{El «prodigio» de Bell Laboratories}{Contexto y red de confianza ciega}
  {\def\usachcifra{\usachtitular\fontsize{22}{24}\selectfont}
  \usachdatos{%
    \begin{minipage}[t]{0.23\linewidth}\raggedright\hyphenpenalty=10000\exhyphenpenalty=10000 \usachdato{8 días}{un paper de alto impacto}\end{minipage}\hfill
    \begin{minipage}[t]{0.23\linewidth}\raggedright\hyphenpenalty=10000\exhyphenpenalty=10000 \usachdato{4 años}{período de publicación (1998–2001)}\end{minipage}\hfill
    \begin{minipage}[t]{0.23\linewidth}\raggedright\hyphenpenalty=10000\exhyphenpenalty=10000 \usachdato{31 años}{candidato al Nobel}\end{minipage}\hfill
    \begin{minipage}[t]{0.23\linewidth}\raggedright\hyphenpenalty=10000\exhyphenpenalty=10000 \usachdato{20}{coautores secundarios}\end{minipage}}}
  \begin{itemize}\usachaire
    \item Bell Labs (EE.~UU.) y Universidad de Konstanz (Alemania): física de
      la materia condensada y nanotecnología.
    \item Transistores monomoleculares y superconductividad en fulerenos, en
      \emph{Science}, \emph{Nature} y APL.
    \item Bertram Batlogg: coautor sénior de fama mundial que nunca presenció
      una medición.
    \item El \emph{efecto halo} de Bell Labs silenció todo escepticismo.
  \end{itemize}
  \note{Schön era la máxima promesa de la nanotecnología: un paper cada ocho
  días y candidato al Nobel a los 31. Batlogg y el prestigio de Bell Labs
  crearon una burbuja de confianza ciega.}
\end{frame}

\begin{frame}{Transgresiones: fabricar, falsificar y borrar}{«On Being a Scientist» y la Declaración de Singapur}
  \begin{columns}[T,onlytextwidth]
    \begin{column}{0.48\textwidth}
      \begin{alertblock}{En «On Being a Scientist» (p.~35)}
        \begin{itemize}
          \item Curvas generadas con fórmulas ideales + ruido artificial.
          \item «Borré los archivos crudos por falta de espacio»; sin bitácoras.
          \item Los chips revolucionarios «se dañaron y los boté a la basura».
          \item Coautores sénior sin validar datos brutos.
        \end{itemize}
      \end{alertblock}
    \end{column}
    \begin{column}{0.48\textwidth}
      \begin{block}{Contravenciones a Singapur}
        \begin{itemize}
          \item Honestidad: adulteración para ventaja profesional.
          \item Resp.~3: conclusiones sin evidencia empírica.
          \item Resp.~4: sin registros reproducibles.
          \item Resp.~6: sin corresponsabilidad de los firmantes.
        \end{itemize}
      \end{block}
    \end{column}
  \end{columns}
  \note{Schön no hacía ciencia experimental: graficaba funciones ideales con
  ruido simulado. Cuando se le pidieron los datos crudos, alegó haberlos
  borrado y haber botado los chips. Quiebre directo de honestidad y de la
  Resp.~4 de Singapur.}
\end{frame}

\begin{frame}{El hallazgo del ruido gemelo}{Revelación y consecuencias históricas}
  \begin{columns}[T,onlytextwidth]
    \begin{column}{0.46\textwidth}
      \textbf{La revelación (2002)}
      \begin{itemize}
        \item Años de réplicas fallidas en Cornell, Harvard, Princeton y Delft.
        \item Sohn (Princeton) y McEuen (Cornell) cotejaron dos papers
          distintos de Schön.
        \item El ruido experimental era \alert{idéntico píxel por píxel}:
          físicamente imposible.
      \end{itemize}
    \end{column}
    \begin{column}{0.50\textwidth}
      \textbf{Comité Beasley y sanciones}
      \begin{itemize}
        \item Fraude confirmado en al menos 16 papers.
        \item \emph{Science} retractó 8 y \emph{Nature} 7 artículos.
        \item Despido inmediato de Bell Labs.
        \item Doctorado revocado por Konstanz (ratificado en 2011).
        \item Protocolos mundiales de conservación de datos crudos.
      \end{itemize}
    \end{column}
  \end{columns}
  \note{No lo detectaron los revisores: lo detectó la comunidad al no poder
  replicar. Dos gráficos de experimentos distintos compartían el mismo ruido
  térmico, copiado píxel a píxel. El Comité Beasley confirmó el fraude y las
  sanciones fueron históricas.}
\end{frame}

% =============================================================================
\section{Caso 2: escándalo ACM}
% =============================================================================

\begin{frame}{ACM: paper mills y mercantilización de la autoría}{El ecosistema afectado}
  {\def\usachcifra{\usachtitular\fontsize{19}{22}\selectfont}
  \usachdatos{%
    \begin{minipage}[t]{0.30\linewidth}\raggedright\hyphenpenalty=10000\exhyphenpenalty=10000 \usachdato{>320}{artículos retractados}\end{minipage}\hfill
    \begin{minipage}[t]{0.30\linewidth}\raggedright\hyphenpenalty=10000\exhyphenpenalty=10000 \usachdato{\$500–3.000}{por firma de autoría}\end{minipage}\hfill
    \begin{minipage}[t]{0.30\linewidth}\raggedright\hyphenpenalty=10000\exhyphenpenalty=10000 \usachdato{48 h}{aprobación sin revisión}\end{minipage}}}
  \begin{itemize}\usachaire
    \item ACM: la sociedad de computación más prestigiosa (Premio Turing);
      actas ICPS afectadas (ICICN, ICISS) en 2021–2022.
    \item \emph{Paper mills}: artículos «llave en mano» con autorías subastadas.
    \item Incentivos perversos: bonos por publicar en Scopus/ACM.
  \end{itemize}
  \note{La ACM sufrió una modalidad nueva: fábricas de artículos que vendían
  autorías por miles de dólares. Ya no es un investigador solitario: es una
  industria ilícita organizada.}
\end{frame}

\begin{frame}{«Tortured phrases» y secuestro del peer review}{Cómo burlaron los filtros antiplagio}
  \begin{table}
    \centering\footnotesize
    \begin{tabularx}{\textwidth}{@{}X X X@{}}
      \toprule
      \rowcolor{usachgris}
      \usachcab{Término real} & \usachcab{Tortured phrase} & \usachcab{Contexto} \\
      \midrule
      Artificial Intelligence & Counterfeit Consciousness & Modelos de ML / IA \\
      Deep Learning & Profound Learning & Redes neuronales \\
      Cloud Computing & Haze Computing & Arquitecturas en la nube \\
      Random Forest & Arbitrary Timberland & Clasificación supervisada \\
      \bottomrule
    \end{tabularx}
    \caption{Ejemplos de parafraseo inverso algorítmico detectados.}
  \end{table}
  \usachfuente{Cabanac, Labbé y Magazinov (2021).}
  \begin{itemize}
    \item \emph{Text spinners} burlan Turnitin/iThenticate por sinónimos inversos.
    \item \emph{Peer-review rings}: revisores ficticios que aprueban en 48 h.
    \item Quiebre de Resp.~8 (arbitraje) y Resp.~6 (autoría comprada).
  \end{itemize}
  \note{Los parafraseadores automáticos no entienden la jerga: «Artificial
  Intelligence» se volvió «Counterfeit Consciousness» y «Random Forest»,
  «Arbitrary Timberland». Igual se aprobaron: los autores controlaban las
  cuentas de revisión.}
\end{frame}

\begin{frame}{Minería de texto y sanciones ACM}{Cómo la propia computación desenmascaró el fraude}
  \begin{columns}[T,onlytextwidth]
    \begin{column}{0.48\textwidth}
      \textbf{Los detectives digitales}
      \begin{itemize}
        \item Cabanac (U.~Toulouse III): \emph{Problematic Paper Screener}.
        \item Labbé y Magazinov: detección de textos sintéticos.
        \item Evidencias publicadas en PubPeer; cobertura de \emph{Nature} y
          Retraction Watch.
      \end{itemize}
    \end{column}
    \begin{column}{0.48\textwidth}
      \textbf{Consecuencias}
      \begin{itemize}
        \item Abril 2022: retractación masiva de \alert{>320 artículos}.
        \item Veto de publicación de 3 a 10 años.
        \item Validó a Beall's List: el modelo predatorio penetró a la ACM.
        \item \emph{Data poisoning}: papers falsos contaminan datasets de IA.
      \end{itemize}
    \end{column}
  \end{columns}
  \note{Cabanac, Labbé y Magazinov programaron un detector de frases
  torturadas; la presión de PubPeer y Retraction Watch forzó a la ACM a
  retractar más de 320 artículos. Y esos papers falsos contaminan los datos
  de entrenamiento de las IA.}
\end{frame}

% =============================================================================
\section{Comparación y lecciones}
% =============================================================================

\begin{frame}{Schön vs. ACM}{La evolución del fraude científico en dos décadas}
  \begin{table}
    \centering\footnotesize
    \begin{tabularx}{\textwidth}{@{}X X X@{}}
      \toprule
      \rowcolor{usachgris}
      \usachcab{Dimensión} & \usachcab{Caso 1: Schön (física)} &
      \usachcab{Caso 2: ACM (informática / IA)} \\
      \midrule
      Naturaleza & Artesanal e individual & Industrializada y sistémica \\
      Mecanismo & Curvas ideales + ruido simulado & \emph{Text spinners} antiplagio \\
      Falla & Confianza ciega en el prestigio & Control editorial tercerizado \\
      Detección & Réplica fallida + ruido idéntico & Minería de texto y PubPeer \\
      Norma rota & FFP + destrucción de bitácoras (Resp.~4) & Autoría comprada (Resp.~6) y arbitraje corrupto (Resp.~8) \\
      \bottomrule
    \end{tabularx}
    \caption{Evolución del fraude: de la oficina individual a la red transnacional.}
  \end{table}
  \usachfuente{elaboración propia a partir de Beasley et al. (2002), Cabanac et al. (2021) y Retraction Watch (2022).}
  \note{En dos décadas el fraude pasó del científico solitario al crimen
  organizado. A Schön lo delató la imposibilidad de replicar; a la ACM, la
  propia computación mediante minería de texto.}
\end{frame}

\usachdestacado{Del fraude artesanal al industrializado, la integridad
científica es el único activo que valida nuestra labor.}

\begin{frame}{Lecciones para el posgrado}{Investigación aplicada en informática}
  \begin{columns}[T,onlytextwidth]
    \begin{column}{0.48\textwidth}
      \begin{block}{Pilares innegociables}
        \begin{itemize}
          \item \textbf{Datos y código abiertos}: decir «perdí los datos» es
            confesar mala praxis.
          \item \textbf{Autoría solidaria}: no existe la autoría honoraria inocua.
          \item \textbf{Peer review como deber moral}: sin bots ni acuerdos recíprocos.
        \end{itemize}
      \end{block}
    \end{column}
    \begin{column}{0.48\textwidth}
      \begin{exampleblock}{Desafíos de la era IA}
        \begin{itemize}
          \item LLM como asistencia, nunca como generador de pseudo-conocimiento.
          \item Vigilancia post-publicación: PubPeer y Retraction Watch.
          \item Superar el \emph{publish or perish}: rigor antes que volumen.
        \end{itemize}
      \end{exampleblock}
    \end{column}
  \end{columns}
  \note{La reproducibilidad y los datos abiertos son un deber ético; la
  autoría implica responsabilidad solidaria; y ante la masificación de la IA
  debemos ser guardianes de la calidad epistemológica de los datos.}
\end{frame}

% --- Fuentes -----------------------------------------------------------------
\begin{frame}{Fuentes y recursos consultados}
  \begin{itemize}\usachaire
    \item National Academy of Sciences. \emph{On Being a Scientist} (3.ª ed.) — caso Schön, p.~35.
    \item World Conferences on Research Integrity. \emph{Declaración de Singapur} (2010).
    \item Retraction Watch Database — reportes de retracción ACM y Bell Labs.
    \item Cabanac, Labbé y Magazinov (2021). \emph{Tortured phrases: a dubious writing style emerging in science}.
    \item Beall's List of Potential Predatory Journals and Publishers.
  \end{itemize}
\end{frame}

\usachcierre[¿Preguntas y comentarios?]{Gracias — Grupo 1 de Cátedra}

\end{document}
